\emph{Översikt}
Övningen är sju uppgifter om interaktiva program: program som frågar den
som kör dem, tål svar som inte går att använda, väljer väg och gör om
något så länge det behövs.  Den börjar i det enklaste --- att
\emph{läsa} ett villkor och räkna ut vad det blir värt --- fortsätter med
ett färdigt program som fungerar men är obehagligt att använda, och
slutar i program som skrivs från början och som ska tåla vad som helst.
De två sista uppgifterna är fördjupning.  Varje uppgift följs av ett
lösningsförslag skrivet som ett program med förklaring, så att den som
inte hann med under övningen kan arbeta vidare på egen hand.

\emph{Lärandemål}
Efter övningen ska studenten kunna:
% Lo04, Lo03
\begin{restatable}{lo}{OvningConditionalsLOBoolean}%
  \label{OvningConditionalsLOBoolean}%
  Skriva villkor med jämförelseoperatorer och de booleska operatorerna
  \mintinline{python}{and}, \mintinline{python}{or} och
  \mintinline{python}{not}, och förutsäga om ett villkor blir sant eller
  falskt.
\end{restatable}
% Lo04
\begin{restatable}{lo}{OvningConditionalsLOBranch}%
  \label{OvningConditionalsLOBranch}%
  Styra programmets flöde med \mintinline{python}{if},
  \mintinline{python}{elif} och \mintinline{python}{else}, och redogöra
  för att exakt en gren exekveras och att programmet fortsätter efteråt.
\end{restatable}
% Lo11
\begin{restatable}{lo}{OvningConditionalsLOTrace}%
  \label{OvningConditionalsLOTrace}%
  Följa exekveringen genom villkor och slingor för hand, ett steg i
  taget, för att förutsäga vad ett program gör och hitta fel i det.
\end{restatable}
% Lo04, Lo01
\begin{restatable}{lo}{OvningConditionalsLOLoop}%
  \label{OvningConditionalsLOLoop}%
  Upprepa en del av programmet tills ett villkor är uppfyllt med
  \mintinline{python}{while}, och avgöra när uppgiften kräver en
  upprepning och när ett val räcker.
\end{restatable}
% Lo08, Lo03
\begin{restatable}{lo}{OvningConditionalsLOConvert}%
  \label{OvningConditionalsLOConvert}%
  Omvandla en inläst sträng till \mintinline{python}{int} eller
  \mintinline{python}{float} med typkonvertering, och avgöra var i
  programmet omvandlingen hör hemma.
\end{restatable}
% Lo04, Lo08
\begin{restatable}{lo}{OvningConditionalsLOCatch}%
  \label{OvningConditionalsLOCatch}%
  Fånga särfall med \mintinline{python}{try} och
  \mintinline{python}{except} --- de specifika sorterna före den
  allmänna --- och redogöra för vad som händer i programmet.
\end{restatable}
% Lo03, Lo04
\begin{restatable}{lo}{OvningConditionalsLORaise}%
  \label{OvningConditionalsLORaise}%
  Låta en egen funktion kasta ett särfall med
  \mintinline{python}{raise} när den inte kan svara vettigt, och låta
  den som anropat funktionen bestämma vad som ska hända.
\end{restatable}
% Lo05, Lo08
\begin{restatable}{lo}{OvningConditionalsLOPrompt}%
  \label{OvningConditionalsLOPrompt}%
  Utforma frågor, utskrifter och felmeddelanden så att den som kör
  programmet förstår vad som förväntas, vad svaret betydde och vad hon
  kan göra när något gick fel.
\end{restatable}
% Lo05, Lo08, Lo01
\begin{restatable}{lo}{OvningConditionalsLORobust}%
  \label{OvningConditionalsLORobust}%
  Utforma en inmatningsdialog som frågar om igen tills svaret går att
  använda, i stället för att upprepa samma fråga i koden.
\end{restatable}
% Lo11
\begin{restatable}{lo}{OvningConditionalsLOReview}%
  \label{OvningConditionalsLOReview}%
  Granska ett färdigt interaktivt program, sätta ord på vad i det som
  brister för den som kör det, och åtgärda bristerna.
\end{restatable}

\emph{Förkunskaper}
Studenten bör ha tagit del av föreläsningen \emph{Inmatning och
felhantering}: att \mintinline{python}{input} ger en sträng, att
strängen måste omvandlas med \mintinline{python}{int} eller
\mintinline{python}{float}, strängmetoderna
\mintinline{python}{strip}, \mintinline{python}{capitalize} och
\mintinline{python}{isdigit}, hur en spårutskrift läses, och hur
\mintinline{python}{try} och \mintinline{python}{except} tar hand om en
omvandling som misslyckas.  Studenten bör också ha tagit del av
föreläsningen \emph{Villkor och styrstrukturer}: jämförelser, de
booleska operatorerna, \mintinline{python}{if}, \mintinline{python}{elif}
och \mintinline{python}{else}, och \mintinline{python}{while}.  Från
föreläsningen \emph{Funktioner} förutsätts egna funktioner med
parametrar, standardvärden och \mintinline{python}{return}, och från
\emph{Algoritmiskt tänkande} principen \foreignlanguage{english}{DRY}.
